\documentclass[11pt,a4paper]{article}
\usepackage[T1]{fontenc}
\usepackage{lmodern}
\usepackage[margin=27mm]{geometry}
\usepackage{amsmath,amssymb,amsthm,booktabs,microtype}
\usepackage[hidelinks]{hyperref}
\newtheorem{theorem}{Theorem}
\newtheorem{lemma}[theorem]{Lemma}
\newtheorem{proposition}[theorem]{Proposition}
\theoremstyle{remark}
\newtheorem{remark}[theorem]{Remark}
\newcommand{\Ostar}{O^{*}}
\newcommand{\E}{\mathbb E}
\newcommand{\Prb}{\mathbb P}
\emergencystretch=2em
\clubpenalty=10000
\widowpenalty=10000
\title{Max 2-SAT by Early Constraint-Guided Enumeration\\
\large A density-independent bound for random signs on almost regular graphs}
\author{}
\date{8 October 2026}
\begin{document}
\maketitle

\begin{abstract}
We prove that the branching search can stop learning after assigning half
the variables and then enumerate only completions permitted by the sum of
its learned constraints. For any fixed simple graph of minimum degree
$d\ge6$ and maximum degree $D$, independently uniform literal signs give
an exact deterministic polynomial-space algorithm with running time
$\Ostar((2e^{-d/(1250D)})^n)$ except on a set of signings of probability
at most $3e^{-nd/(1250D)}$. In particular, every regular interaction graph
gives a base below $1.999$, independently of its degree. The proof counts
all generated completions, including exceptional prefixes. It does not
require a conflict at the switching depth or a spectral hypothesis.
The input guarantee is probabilistic over signs, not a worst-case bound
for every signing of an almost regular graph. Constants are deliberately
conservative, and no literature novelty claim is made.
\end{abstract}

\section{The input model and result}

Fix a simple graph $G$ on $n$ labeled vertices. Each edge $uv$ contributes
one clause $(\ell_u\vee\ell_v)$, whose two literal signs are independent
uniform elements of $\{-1,+1\}$. All signs on distinct incidences are
independent. There are no initial unit clauses. The graph and a vertex
order are fixed before the signs are drawn. The graph itself can be
arbitrary, including a complete graph; it need not be randomly generated.

Write $d=\min_v\deg(v)$ and $D=\max_v\deg(v)$. These denote the minimum
and maximum graph degrees, not the mean squared row norm used in the
explicit-class note. All logarithms without a subscript are natural.

\begin{theorem}[Early enumeration]\label{thm:main}
There is a deterministic algorithm that is exact on every instance in
this model and uses polynomial space. If $d\ge6$, then with probability
at least $1-3e^{-nd/(1250D)}$ over the signs, its running time is
\[
 \Ostar\left((2e^{-d/(1250D)})^n\right).
\]
The algorithm branches in the fixed order through depth
$T=\lceil n/2\rceil$, learning the usual constraint in each nonpreferred
child. It then freezes those constraints and enumerates completions of
their summed inequality. No contradiction at depth $T$ is required.
\end{theorem}

For $d$-regular graphs, the base is
$2e^{-1/1250}\simeq1.99840064<1.999$.
More generally, $D/d\le R$ for any fixed $R$ gives a fixed base below two.
Degrees may grow with $n$, up to $n-1$; in the regular case the clause
density is $d/2$. The expectation of the running time satisfies the
stronger bound $\Ostar((2e^{-d/(625D)})^n)$.
The algorithm does not need to recognize a typical signing or know $R$.

This is a precisely specified random-input result. It is not a theorem
for all signings of regular graphs, and it does not supply a combinatorial
membership test for the typical signings. The earlier notes already
obtain density-independent bounds in this random-sign model by waiting
for conflicts~\cite{Key}. The point here is to justify an earlier switch
to enumeration, using a different quantity in the analysis.

\section{Algorithm, exactness, and cost}

Use Boolean values $\sigma_v\in\{0,1\}$. At step $i$, let $\Delta_i$ be
the number of positive occurrences minus negative occurrences of $x_i$
in residual 2-clauses, plus twice the corresponding difference in
residual unit clauses, after assigning $1,\ldots,i-1$.
Let $D_i^+$ and $D_i^-$ be
the numbers of neighbors later and earlier in the fixed order.
The preferred literal is positive if $\Delta_i\ge0$, and negative
otherwise. In the nonpreferred child learn
\begin{equation}
 L_i\le k_i^0,
 \qquad k_i^0=\frac{D_i^+-|\Delta_i|}{2}-1,          \label{eq:learn}
\end{equation}
Here $L_i$ counts one bad value per later neighbor: in a clause containing
the preferred literal of $x_i$, the neighbor value is bad if its literal
is false; in a clause containing the other literal of $x_i$, it is bad
if its literal is true. This is the previous branching
rule~\cite{Plan,Frontier}. The right-hand side is integral because
$\Delta_i$ and $D_i^+$ have the same parity. Negative allowances can be
cut immediately.

For clarity, the upper-bound analysis uses the \emph{uncut process}:
generate the constraints for every prefix, even if a negative allowance
or an earlier test would already have killed it. Early cuts can only
reduce the actual work. No new constraints are created after step $T$.

At a prefix $\alpha\in\{0,1\}^T$, substitute its values into every
created constraint. If the remaining constraints are $L_c(y)\le k_c$,
put $K=\sum_c k_c$. For each free variable $w$, let $m_w^b$ count the
occurrences made true by $y_w=b$. Orient its cheaper value as $z_w=0$ and
define
\begin{equation}
 c_w=|m_w^1-m_w^0|,\qquad
 B=K-\sum_w\min\{m_w^0,m_w^1\}.                    \label{eq:budget}
\end{equation}
The sum of the constraints is equivalent to
$\sum_w c_wz_w\le B$. Let $N_F(\alpha)$ be its number of Boolean
solutions, including zero when $B<0$.

Generate those solutions by depth-first search, pruning whenever the
current nonnegative cost exceeds $B$. Evaluate the original Max 2-SAT
score on each generated complete assignment. Testing every individual
learned constraint as well is optional: the summed condition is a
necessary relaxation, and extra candidates are still valid assignments
of the original soft optimization problem.

\begin{lemma}[Exactness and enumeration cost]\label{lem:cost}
The algorithm preserves an optimum and has time
\begin{equation}
 \Ostar\left(2^T+W(F)\right),\qquad
 W(F)=\sum_{\alpha\in\{0,1\}^T}N_F(\alpha),        \label{eq:cost}
\end{equation}
using polynomial space.
\end{lemma}
\begin{proof}
For a fixed prefix and suffix, the score of the preferred value minus
that of the nonpreferred value at decision $i$ is
$L_i-(D_i^+-|\Delta_i|)/2$. Consequently, violating~\eqref{eq:learn}
means that switching to the preferred value is no worse.
Choose the first optimal assignment in the full search-tree order that
visits preferred children first. If it violates any learned constraint
on its path, flipping that decision gives an equally good assignment
earlier in that order, a contradiction. Thus some optimum satisfies
all its learned constraints, and hence their sum.

For a fixed prefix with $N=N_F(\alpha)>0$, every retained generator
prefix extends to a permitted completion by taking the cheaper value
on every remaining variable. At each generator depth there are at most
$N$ retained nodes, and there are at most $n-T+1$ depths. Rejected
children add only a constant factor. If $N=0$, the summed inequality
has $B<0$ and is rejected immediately. Thus each prefix costs
$\operatorname{poly}(n)(N+1)$, yielding~\eqref{eq:cost}.
Polynomial-size constraints and a depth-first traversal give polynomial
space. Retaining or recomputing a polynomial number of integer counters
suffices; a table over exponentially many prefixes is unnecessary.
\end{proof}

\section{Counting completions instead of waiting for a conflict}

Let $\sigma\in\{0,1\}^n$ be a uniformly random complete assignment,
independent of the formula. Run its first $T$ decisions through the
uncut process, and let $\mathcal C_T$ denote the created constraints.
For $c$ created at $i$, let $L_c(\sigma)$ count all its bad values on
the later neighbors, before any substitutions. Define
\begin{align}
 S_T&=\sum_{c\in\mathcal C_T}\left(\frac{|\Delta_c|}{2}+1\right),
                                                        \label{eq:S}\\
 Q_T^{\rm all}&=\sum_{c\in\mathcal C_T}
                 \left(L_c(\sigma)-\frac{D_c^+}{2}\right).
                                                        \label{eq:Q}
\end{align}
The superscript distinguishes this from the earlier $Q_T$, which counted
only support variables assigned by depth $T$. Here all support variables
are evaluated, but only constraints created by time $T$ are used.

\begin{lemma}[Candidate-mass identity]\label{lem:mass}
For every complete assignment in the uncut process,
\[
 \sum_{c\in\mathcal C_T}(L_c(\sigma)-k_c^0)
   =S_T+Q_T^{\rm all}.
\]
Consequently,
\begin{equation}
 \frac{\E_F W(F)}{2^n}
 =\Prb_{F,\sigma}\left[S_T+Q_T^{\rm all}\le0\right].
                                                        \label{eq:mass}
\end{equation}
\end{lemma}
\begin{proof}
Substitute~\eqref{eq:learn} and add. For a fixed $F$, the number of full
assignments satisfying the summed condition belonging to their own
prefix is precisely $W(F)$. Divide by $2^n$ and average over $F$.
\end{proof}

This identity replaces the requirement that an entire prefix be
infeasible. A prefix can have many candidates and still contribute
little to the total in~\eqref{eq:cost}. The proof averages their actual
counts, so unusually expensive prefixes are included.

\section{Two lower-tail estimates}

Let $H_i$ reveal the signs of vertices $1,\ldots,i$ in all their incident
clauses, together with $\sigma_1,\ldots,\sigma_i$. Future endpoint signs
are not revealed. The uncut process uses this filtration; we never
condition it on surviving an earlier cut. A fixed order is important
for this exposure argument.

\subsection{Accumulated shortfall}

At step $i$, let $r_i$ be the number of earlier neighbors whose literal
in their clause with $i$ is false. This is determined by $H_{i-1}$.
Conditional on $H_{i-1}$, the gap $\Delta_i$ is a sum of independent
uniform signs: $D_i^+$ signs of weight one and $r_i$ signs of weight two.
Put $Y_i=D_i^++4r_i$.

We use an elementary moment bound rather than the optimal Khintchine
constant. For a Rademacher sum $Z$ with variance $v$, expansion gives
$\E Z^4\le3v^2$. H\"older's inequality gives
$\E Z^2\le(\E|Z|)^{2/3}(\E Z^4)^{1/3}$, hence
\[
 \E|Z|\ge\sqrt{v/3}.
\]
The bound is also valid when $v=0$.

Let $X_i$ be the contribution of decision $i$ to $S_T$.
Given its gap, the decision is nonpreferred with probability $1/2$.
Thus
\begin{align*}
 \E[X_i\mid H_{i-1}]&\ge
       \frac{\sqrt{Y_i}}{4\sqrt3}+\frac12=:\mu_i,\\
 \E[X_i^2\mid H_{i-1}]&\le
       \frac{Y_i}{8}+\frac{\sqrt{Y_i}}2+\frac12
       \le\frac D2+\sqrt D+\frac12\le D,
\end{align*}
where $Y_i\le4D$ and $D\ge6$.
For $X_i\ge0$, the inequality
$e^{-\lambda X_i}\le1-\lambda X_i+\lambda^2X_i^2/2$ implies a
conditional exponential-moment bound. Multiplying its compensated
versions and taking $\lambda=a/(TD)$ gives
\begin{equation}
 \Prb\left[S_T<\sum_{i\le T}\mu_i-a\right]
 \le e^{-a^2/(2TD)}.                                \label{eq:Stail}
\end{equation}

It remains to lower-bound the sum of the predictable means.
The indicators of earlier endpoint literals being false are independent
fair coins: conditional on all of $\sigma$, they use distinct independent
incidence signs, and their joint distribution does not depend on
$\sigma$. Thus $r_i\sim\operatorname{Bin}(D_i^-,1/2)$, with
\[
 \E Y_i=D_i^++2D_i^-\ge d,\qquad
 \operatorname{Var}(Y_i)=4D_i^-\le2\E Y_i.
\]
For a nonnegative random variable $Y$ of mean $\mu>0$,
\[
 \E\sqrt Y\ge\sqrt\mu-
                  \frac{\operatorname{Var}(Y)}{2\mu^{3/2}}.
\]
This follows by averaging the pointwise inequality
$\sqrt z\ge1+(z-1)/2-(z-1)^2/2$ for $z\ge0$; after writing
$z=t^2$, the difference factors as $t(t-1)^2(t+2)/2\ge0$.
It follows that $\E\sqrt{Y_i}\ge\sqrt{D_i^++2D_i^-}-1$.
The positive $1/2$ in $\mu_i$ pays for this subtraction, so
\[
 \E\sum_{i\le T}\mu_i\ge\frac{T\sqrt d}{4\sqrt3}=:\mu_*.
\]

Let $e_P$ count edges internal to the first $T$ vertices. Changing one
of its earlier-endpoint coins changes $\sum_i\mu_i$ by at most
$1/(2\sqrt3)$. The bounded-differences inequality therefore gives
\begin{equation}
 \Prb\left[\sum_{i\le T}\mu_i<\mu_*-b\right]
 \le e^{-24b^2/e_P}
 \le e^{-48b^2/(TD)}.                               \label{eq:mutail}
\end{equation}
If $e_P=0$, the sum is deterministic and this failure probability is zero.
Combining~\eqref{eq:Stail}--\eqref{eq:mutail} controls $S_T$.

\subsection{All support contributions}

Reveal variables through $n$, but create no more constraints after $T$.
When vertex $w$ is revealed, its contribution to $Q_T^{\rm all}$ is
\[
 \sum_{\substack{i<w,\ i\le T\\ i\text{ created a constraint}}}
       \left([\sigma_w\text{ is bad in the constraint of }i]-\frac12\right),
\]
with only adjacent $i$ included. Conditional on $H_{w-1}$, the creator
set is fixed. The bad values use distinct unseen signs of $w$ and are
independent fair coins. Multiplication by the common independent value
of $\sigma_w$ preserves their joint uniform distribution.

Let $a_w$ count neighbors $i<w$ with $i\le T$, whether or not they
created a constraint. The conditional moment generating function of
the negative contribution is at most
$\cosh(\lambda/2)^{a_w}\le e^{\lambda^2a_w/8}$.
If $e_\times$ counts edges from the prefix to its complement, then
\[
 \sum_w a_w=e_P+e_\times\le TD.
\]
Multiplying the conditional bounds and optimizing $\lambda$ proves
\begin{equation}
 \Prb[Q_T^{\rm all}<-c]
 \le e^{-2c^2/(e_P+e_\times)}
 \le e^{-2c^2/(TD)}.                               \label{eq:Qtail}
\end{equation}
No independence between $S_T$ and $Q_T^{\rm all}$ is asserted or needed.

\section{Completing the running-time proof}

Choose
\[
 a=\frac{2}{25}T\sqrt d,\qquad
 b=\frac1{100}T\sqrt d,\qquad
 c=\frac1{20}T\sqrt d.
\]
Their sum is $(7/50)T\sqrt d$, strictly less than
$\mu_*=T\sqrt d/(4\sqrt3)$. Outside the three exceptional events,
$S_T+Q_T^{\rm all}>0$. The exponents in
\eqref{eq:Stail}, \eqref{eq:mutail}, and \eqref{eq:Qtail} are at least,
respectively,
\[
 \frac{2Td}{625D},\qquad
 \frac{3Td}{625D},\qquad
 \frac{Td}{200D}.
\]
Hence Lemma~\ref{lem:mass} yields the following bound for \emph{every}
$1\le T\le n$:
\begin{equation}
 \E_F W(F)\le3\cdot2^n e^{-2Td/(625D)}.             \label{eq:EW}
\end{equation}
By Markov's inequality,
\begin{equation}
 \Prb_F\left[W(F)>2^n e^{-Td/(625D)}\right]
 \le3e^{-Td/(625D)}.                               \label{eq:Markov}
\end{equation}
For $T=\lceil n/2\rceil$, the right-hand side is at most
$3e^{-nd/(1250D)}$, and outside this exceptional set
\[
 W(F)\le(2e^{-d/(1250D)})^n.
\]
The $2^T$ term in~\eqref{eq:cost} is smaller up to a constant factor.
This proves Theorem~\ref{thm:main}. Applying~\eqref{eq:EW} directly to
the cost bound proves the stated expectation bound.\hfill$\square$

\paragraph{Earlier switching.}
For any fixed $\tau\in(0,1)$, taking $T=\lceil\tau n\rceil$ instead
gives, with failure probability at most $3e^{-\tau nd/(625D)}$,
\[
 \Ostar\left(2^{\tau n}+2^n e^{-\tau nd/(625D)}\right).
\]
Thus every fixed positive switching fraction gives a fixed saving for
bounded $D/d$. Halfway is a simple concrete choice, not a required
threshold and not an optimized one.

\section{Interpretation and limitations}

The conflict-only analysis bounds the number of prefixes with any
remaining work. This argument bounds their \emph{total remaining work}.
It pays for a prefix whose summed constraint is satisfiable but admits
few assignments, and handles expensive prefixes through an expectation
bound followed by Markov's inequality. No typical-prefix or median-cost
assumption is substituted for that calculation.

The algorithm can deliberately do no conflict checks before $T$ and
still satisfy the theorem. Earlier checks can be added to the same fixed
search without increasing its cost. Adaptive ordering during the first
phase is not covered: inspecting future signs to choose a variable can
invalidate the late-exposure argument. The final budget generator may
use any variable order because its retained prefixes extend to solutions
of the same frozen inequality. Further learning would require preserving
that inherited inequality or another valid cost certificate.

The result establishes the user's proposed early-enumeration mechanism
for a well-defined random-sign input model, with an elementary proof.
It neither proves the general growing-density conjecture nor resolves
which deterministic interaction patterns guarantee the same bound.
The constants were chosen for a transparent argument, not to improve
the numerical base in the previous random-sign theorem.

\section{Computational checks}

The companion program \texttt{explore\_early\_enumeration.py} checks the
algorithm on small signed regular graphs, including degrees below the
theorem's minimum to test exactness separately from the running-time
guarantee. It compares the optimum with exhaustive search, compares
budget enumeration and exact dynamic-program counts with direct counts,
and verifies the identity of Lemma~\ref{lem:mass} on every complete
assignment. All 44 tested instances on 6, 8, 10, and 12 variables
preserved the optimum; 1,408 prefix counts and 65,024 full-assignment
identities passed. Its results are recorded in
\texttt{early\_enumeration\_checks.json}.
These finite checks validate the implementation and algebra; the
probabilistic running-time guarantee follows from the proof above.

\begin{thebibliography}{9}
\small\raggedright
\bibitem{Frontier}
\emph{Exploring the Branching Frontier for Max 2-SAT}, working note,
7 October 2026. Section 3: a joint measure from exact weighted-sum counts.
Local source: \texttt{max2sat\_branching\_frontier.tex}.
\bibitem{Plan}
\emph{A Simplified Plan for Max 2-SAT: One Search, One Count, One Key Lemma},
working note, 8 October 2026. Local source:
\texttt{max2sat\_plan\_key\_lemma.tex}.
\bibitem{Key}
\emph{The Key Lemma for Random Signs}, working note, 8 October 2026.
Local source: \texttt{max2sat\_key\_lemma\_proof.tex}.
The shortfall and late-exposure arguments are adapted here, with an
elementary fourth-moment constant and a completion-count objective.
\end{thebibliography}
\end{document}
